% TP2 — resumen del enunciado y del dataset Bank Marketing (skill /dossier). Para el grupo: qué pide
% el TP2, qué trae el dataset y qué hay que decidir antes del envío del 06/10.
% Perfil: breve — texto normal — imágenes más — lector equipo — formato digital
% Medir:  python3 ~/.claude/skills/dossier/scripts/medir.py resumen-tp2.tex --paginas 6 --paginas-texto 1.8 --densa 550 --visuales-min 9
% Regenerar:  python3 graficos.py && latexmk -lualatex resumen-tp2.tex
\documentclass[10pt,a4paper]{article}
\usepackage[spanish,es-noshorthands]{babel}
\usepackage{dossier}
\documento{72.75 · TP2 · Clasificación supervisada}{TP2: quién contrata un plazo fijo}{2C 2026}

\begin{document}

% ============================ PÁGINA 1 ============================
\portada{72.75 Aprendizaje Automático · TP2 · 2C 2026}
  {TP2: predecir quién contrata un plazo fijo}
  {Qué pide el enunciado, qué trae el dataset Bank Marketing y qué hay que decidir antes del envío}
  {Documento de lectura para el grupo, 25/09/2026. Fuentes en gris: \textit{TP2} = enunciado;
  \textit{CSV} = bank-additional-full.csv, cálculo propio; \textit{names} = bank-additional-names.txt.}

\hitos[0.07]{0/Enunciado/25-09/hecho,0.25/Clase KNN?/30-09/pendiente,0.5/Envío/06-10/pendiente,0.75/1.ª defensa/07-10/pendiente,1/2.ª defensa/14-10/pendiente}
{\footnotesize\color{apagado} Fechas: \fuente{TP2, p. 1 · cronograma de la Clase 1}. KNN el 30/09 es \chip{inferencia}; el envío es el 06/10 si el grupo defiende el 07/10.}

\begin{cifras}[4]
  \cifra{Llamadas}{41.188}{20 predictoras + $y$ \fuente{CSV}}
  \cifra{Contratan}{11,27\%}{4.640; decir siempre «no» acierta 88,73\% \fuente{CSV}}
  \cifra{Clasificadores}{4}{Naive Bayes, SVM, KNN, RF \fuente{TP2, p. 2}}
  \cifra{Defensa}{10 + 10}{minutos; envío 24 h antes \fuente{TP2, p. 1}}
\end{cifras}

\resumen
\textbf{Qué pide.} Predecir si un cliente contrata un depósito a plazo fijo ($y$ = yes/no) con
Naive Bayes, SVM, KNN y Random Forest en scikit-learn \fuente{TP2, pp. 1–2}. Se evalúa con
\textit{k-fold} sólo sobre entrenamiento, dos métricas justificadas y curvas de validación para al
menos un hiperparámetro de SVM, KNN y RF; el modelo final estima su desempeño en datos nuevos. El
EDA puede ser básico, pero sus observaciones tienen que tener consecuencias en el pipeline.

\textbf{Qué trae el dato.} Una llamada de cada nueve termina en «sí», así que la exactitud sola
no sirve. Hay tres trampas: \texttt{duration}
mira el futuro, hay faltantes disfrazados (\texttt{unknown}, \texttt{pdays = 999}) y el archivo está
ordenado por fecha: la tasa de «sí» pasa de 4,84\% (2008) a 52,14\% (2010, año inferido) \fuente{names §4 · CSV}.

\textbf{Qué sigue.} Seis decisiones antes del envío (\ver{sec:pasos}). Tres conviene consultarlas
con la cátedra: \texttt{duration}, partir al azar o por fecha, y si $F_1$ está entre las métricas vistas.

\ojo[Antes de entrenar]{\texttt{duration} no existe antes de llamar. Los autores piden descartarla
para un modelo realista \fuente{names §7}, y sola separa las clases con AUC 0,818 \fuente{CSV}. Un
modelo que la usa parece excelente y no sirve.}

\indice

\begin{bloque}[La consigna, en el orden del enunciado. Azul oscuro: hecho; celeste: en curso \fuente{TP2, pp. 1–2}.]
\proceso[1]{Dataset/0/hecho,EDA con decisiones/1/actual,Clasificar con CV/2/futuro,Curvas de validación/3/futuro,Modelo final/4/futuro,Conclusiones/5/futuro}
\end{bloque}

% ============================ PÁGINA 2 ============================
\newpage
\section{El enunciado evalúa cómo se valida, más que qué modelo gana}\label{sec:enunciado}

El objetivo declarado es «aprender a entrenar y evaluar modelos utilizando validación cruzada»
\fuente{TP2, p. 1}: el test se separa primero y toda transformación aprendida se ajusta sólo con
entrenamiento.

\begin{bloque}[Nada pasa del carril del test al de entrenamiento. Separar antes de transformar lo pide el enunciado \fuente{TP2, p. 1}; usar el test una vez y reentrenar con todo el train es \chip{inferencia}.]
\begin{tikzpicture}[node distance=5mm]
  \carril{0}{1.95}{ENTRENAMIENTO}
  \carril{-1.8}{1.5}{TEST}
  \node[cajag,fill=white,text width=2.25cm] (csv) at (1.95,0.08) {\textbf{CSV}\\41.188 filas \fuente{CSV}};
  \node[cajag,fill=white,text width=2.75cm,right=11mm of csv,yshift=0.9cm] (kf) {\textbf{k-fold sobre train}\\pipeline ajustado en cada pliegue};
  \node[cajag,fill=white,text width=3.1cm,right=of kf] (cv) {\textbf{Curvas de validación}\\SVM, KNN y RF: uno o más hiperparámetros};
  \node[cajag,fill=white,text width=2.8cm,right=of cv] (mf) {\textbf{Modelo final}\\reentrenado con todo el train};
  \node[cajag,fill=white,text width=2.35cm] (te) at (5.85,-1.05) {\textbf{Test apartado}\\antes de transformar};
  \node[cajaa,text width=2.35cm] (ev) at (15.95,-1.05) {\textbf{Test, una vez}\\desempeño esperado};
  \draw[rel] (csv.east) -- ++(0.4,0) |- (kf.west);
  \draw[rel] (csv.east) -- ++(0.4,0) |- (te.west);
  \draw[rel] (kf) -- (cv); \draw[rel] (cv) -- (mf);
  \draw[rel] (mf.east) -| (ev.north);
  \draw[rel,dashed] (te.east) -- node[relacion,above] {no se toca hasta el final} (ev.west);
\end{tikzpicture}
\end{bloque}

\begin{tabularx}{\linewidth}{@{}P{2.4cm}P{4.2cm}L@{}}
\toprule
\enc{Modelo} & \enc{Qué pide evaluar} & \enc{Qué mirar en este dataset} \\
\midrule
\textbf{Naive Bayes} & ningún hiperparámetro \fuente{TP2, p. 2} & el enunciado no dice qué variante; hay 10 numéricas y 10 categóricas \fuente{CSV} \\
\textbf{SVM} & C o kernel \fuente{TP2, p. 2} & escalar es obligatorio; con 41.188 filas es el más caro de entrenar \chip{inferencia} \\
\textbf{KNN} & número de vecinos, ponderación \fuente{TP2, p. 2} & escalar; con 11,27\% de «sí», un $k$ grande vota «no» casi siempre \chip{inferencia} \\
\textbf{Random Forest} & número de árboles, profundidad máxima \fuente{TP2, p. 2} & no necesita escalar; con \texttt{duration} adentro, los cortes la eligen a ella \chip{inferencia} \\
\bottomrule
\end{tabularx}
\medskip

\textbf{Lo que el enunciado no dice, y queda a criterio del grupo:}
\begin{itemize}
\item \textbf{Qué Naive Bayes} para datos mixtos.
\item \textbf{Cuántos pliegues y si se estratifica.}
\item \textbf{Si se trata el desbalance.} Pide mirar «el balance de clases» \fuente{TP2, p. 1}, no corregirlo.
\item \textbf{Qué variable filtra.} Advierte sobre la fuga sin nombrar ninguna \fuente{TP2, p. 1}.
\item \textbf{El tamaño del test y el umbral de decisión.}
\end{itemize}
Del TP1 arrastra tres consejos \fuente{TP2, pp. 2–3}: decir qué datos se usan en cada etapa, preferir
gráficos a tablas y no pasar de 10 minutos, con los números en las slides. El objetivo dice
«regresión»: copia del TP1 \fuente{TP1, p. 1}.

% ============================ PÁGINA 3 ============================
\newpage
\section{El dataset: 20 predictoras en cuatro grupos y una clase de cada nueve}\label{sec:dataset}

Cada fila es un contacto telefónico de un banco portugués entre mayo de 2008 y noviembre de 2010
\fuente{names §4}. Sólo uno de los cuatro grupos describe la llamada misma.

\begin{bloque}[Los grupos de names §7. \texttt{duration} es del último contacto, pero se conoce recién al colgar: va tras la línea.]
\begin{tikzpicture}
  \node[grupo] (a) at (0,0)     {\grupo[3.2cm]{3.1cm}{Cliente · 7}{age, job, marital, education, default, housing, loan}};
  \node[grupo] (c) at (3.65,0)  {\grupo[3.2cm]{2.6cm}{Otros atributos · 4}{campaign, pdays, previous, poutcome}};
  \node[grupo] (d) at (6.8,0)   {\grupo[3.2cm]{3.2cm}{Contexto económico · 5}{emp.var.rate, cons.price.idx, cons.conf.idx, euribor3m, nr.employed}};
  \node[grupo] (b) at (10.55,0) {\grupo[3.2cm]{2.75cm}{Último contacto · 3}{contact, month, day\_of\_week}};
  \node[grupo] (e) at (14.45,0) {\grupo[3.2cm]{2.3cm}{Al colgar · 1}{!duration}};
  \draw[acento2,dashed,line width=0.9pt] (14.1,0.3) -- (14.1,-3.85);
  \node[font=\scriptsize\bfseries,text=acento2,anchor=south east] at (14.0,0.02) {momento de predecir};
  \coordinate (bus) at (0,-3.95);
  \foreach \g in {a,c,d,b,e} \draw[bus] (\g.south) -- (\g.south |- bus);
  \draw[bus] (a.south |- bus) -- (e.south |- bus);
  \coordinate (medio) at ($(a.south |- bus)!0.5!(e.south |- bus)$);
  \node[concepto central,minimum width=5.2cm] (y) at ($(medio)+(0,-0.75)$) {$y$: ¿contrató el depósito? (yes/no)};
  \draw[rel] (medio) -- (y);
\end{tikzpicture}
\end{bloque}

\noindent\begin{minipage}[t]{0.48\linewidth}\vspace{0pt}
\figura{fig/f01_balance.pdf}{Decir siempre «no» acierta el 88,73\%: es la línea de base de cualquier exactitud \fuente{CSV}.}
\end{minipage}\hfill
\begin{minipage}[t]{0.48\linewidth}\vspace{0pt}
\figura{fig/f02_unknown.pdf}{Ningún NaN: los faltantes vienen como la categoría \texttt{unknown} \fuente{names §8 · CSV}.}
\end{minipage}

\begin{itemize}
\item \textbf{10 numéricas y 10 categóricas.} Las categóricas suman 53 niveles: \textit{one-hot} sin descartar da 53 columnas \fuente{CSV}.
\item \textbf{Un cuarto de las filas tiene algún \texttt{unknown}} (25,98\%); en \texttt{default}, 8.597 contra sólo 3 «yes» \fuente{CSV}.
\item \textbf{Categorías raras.} \texttt{default = yes} (3 filas, ninguna con «sí») e \texttt{illiterate} (18): sin Laplace, Naive Bayes puede darles probabilidad cero \fuente{CSV}.
\item \textbf{12 filas duplicadas exactas}: al partir, una copia puede quedar en cada lado \fuente{CSV}.
\end{itemize}

\figura{fig/f09_categorias.pdf}{Porcentaje de «sí» por nivel; la línea punteada es el 11,27\% global y el gris marca categorías con menos del 1\% de las filas. \texttt{month} lleva la época adentro: marzo, abril y septiembre no tienen filas de 2008, y octubre y diciembre casi ninguna \fuente{CSV}.}


\subsection{\texttt{pdays = 999} no siempre significa «nunca contactado»}
\noindent\begin{minipage}[t]{0.52\linewidth}\vspace{0pt}
\figura{fig/f04_pdays.pdf}{El 96,32\% de las filas tiene 999; 4.110 de ellas sí tuvieron contacto previo \fuente{CSV}.}
\end{minipage}\hfill
\begin{minipage}[t]{0.44\linewidth}\vspace{4pt}
El diccionario define el 999 como «no contactado antes» \fuente{names §7}, pero 4.110 filas con 999
tienen \texttt{previous}~$\geq$~1 \fuente{CSV}. El indicador fiable es \texttt{previous = 0}. Como
número, el 999 aplasta las distancias de KNN y SVM \chip{inferencia}: mejor una indicadora más los
días reales \chip{propuesta}.
\end{minipage}

% ============================ PÁGINA 4 ============================
\newpage
\section{Dos trampas que cambian lo que mide la evaluación}\label{sec:trampas}

\subsection{\texttt{duration} predice porque mira el futuro}
\noindent\begin{minipage}[t]{0.52\linewidth}\vspace{0pt}
\figura{fig/f03_duration.pdf}{Ninguna llamada de menos de 30 segundos terminó en «sí»; las 4 de 0 segundos, en «no» \fuente{CSV}.}
\end{minipage}\hfill
\begin{minipage}[t]{0.44\linewidth}\vspace{4pt}
La duración se conoce cuando la llamada termina, que es cuando ya se sabe la respuesta. Encaja con
la advertencia del enunciado sobre información que no está «antes de realizar la llamada»
\fuente{TP2, p. 1} \chip{inferencia}. Entre las decisiones del EDA, el enunciado nombra «la exclusión
de alguna variable» y comparar el modelo «con y sin ciertas variables problemáticas» \fuente{TP2, p. 1}.
\end{minipage}

\subsection{El archivo está ordenado por fecha, y la tasa de «sí» cambia con ella}
\figura{fig/f06_orden_temporal.pdf}{Sin barajar, cada pliegue es otra época: de 3,1\% a 30,8\% de «sí» \fuente{CSV}.}

\noindent\begin{minipage}[t]{0.3\linewidth}\vspace{0pt}
\figura{fig/f05_por_anio.pdf}{El año no está en el archivo: se infiere del ciclo de meses \fuente{names §4 · CSV}.}
\end{minipage}\hfill
\begin{minipage}[t]{0.66\linewidth}\vspace{6pt}
\texttt{train\_test\_split} baraja por defecto (documentación de scikit-learn) y los pliegues dejan
de ser épocas. Pero la estimación mide el desempeño dentro de las mismas épocas; partir por fecha se
parece más a usar el modelo en campañas futuras \chip{inferencia}. La cátedra no dijo cuál espera.
\end{minipage}

% ============================ PÁGINA 5 ============================
\newpage
\section{Qué le hace este dataset a cada uno de los cuatro modelos}\label{sec:modelos}

El punto 5 pide mencionar si los datos se ajustan a los supuestos de los modelos, con cuatro
ejemplos: variables correlacionadas, distribuciones no gaussianas, diferencias de escala y relaciones
no lineales \fuente{TP2, p. 2}. Los cuatro están en el dataset.

\noindent\begin{minipage}[t]{0.58\linewidth}\vspace{0pt}
\figura{fig/f07_correlacion.pdf}{Recuadro: tres indicadores macro con $r$ entre 0,91 y 0,97. Para Naive Bayes, la misma evidencia contada tres veces \fuente{CSV}.}
\end{minipage}\hfill
\begin{minipage}[t]{0.39\linewidth}\vspace{0pt}
\figura{fig/f08_varianza.pdf}{Sin escalar, \texttt{duration} y \texttt{pdays} son el 95\% de la varianza: una distancia de KNN o un núcleo de SVM mira casi sólo esas dos \fuente{CSV} \chip{inferencia}.}
\medskip
\figura{fig/f10_edad.pdf}{El «sí» sube en los dos extremos de edad: una relación no lineal \fuente{CSV}.}
\end{minipage}

\begin{tabularx}{\linewidth}{@{}P{2.4cm}L@{}}
\toprule
\enc{Modelo} & \enc{Qué le pasa con estos datos} \\
\midrule
\textbf{Naive Bayes} & el bloque macro rompe la independencia condicional; las numéricas son asimétricas y \texttt{pdays} es bimodal, mala forma para una normal por clase \chip{inferencia} \\
\textbf{SVM} & sin escalar, la frontera la deciden \texttt{duration} y \texttt{pdays} \chip{inferencia}; los desvíos van de 0,5 a 259 \fuente{CSV} \\
\textbf{KNN} & la misma escala, más 53 columnas \textit{one-hot} \fuente{CSV} mezcladas con numéricas en una sola distancia \\
\textbf{Random Forest} & inmune a la escala; reparte los cortes entre tres variables macro casi iguales y diluye su importancia \chip{inferencia} \\
\bottomrule
\end{tabularx}

% ============================ PÁGINA 6 ============================
\newpage
\section{Próximos pasos: seis decisiones antes del envío}\label{sec:pasos}

Las fechas intermedias son propuesta. La fija es el envío, 24 h antes de la clase en que se
presenta \fuente{TP2, p. 1}: 06/10 si el grupo defiende el 07/10, 13/10 si defiende el 14/10.

\begin{tabularx}{\linewidth}{@{}LP{3.3cm}P{1.7cm}@{}}
\toprule
\enc{Qué} & \enc{Quién} & \enc{Para} \\
\midrule
\textbf{\texttt{duration}:} excluirla o comparar el modelo con y sin ella & grupo; consultar a la cátedra & 30/09 \chip{propuesta} \\
\textbf{Partición:} barajar y estratificar, o partir por fecha & grupo; consultar a la cátedra & 30/09 \chip{propuesta} \\
\textbf{Faltantes:} \texttt{unknown} como categoría; indicadora \texttt{previous > 0} en lugar del 999 & grupo & 30/09 \chip{propuesta} \\
\textbf{Métricas:} dos que no premien decir siempre «no»; ¿$F_1$ está entre «las vistas en clase»? & grupo; consultar a la cátedra & 02/10 \chip{propuesta} \\
\textbf{Naive Bayes:} qué variante para 10 numéricas y 10 categóricas & grupo & 02/10 \chip{propuesta} \\
\textbf{Escala y colinealidad:} escalar dentro del pipeline para SVM y KNN; decidir si queda una sola variable macro & grupo & 02/10 \chip{propuesta} \\
\textbf{Enviar} presentación y código & grupo & 06/10 o 13/10 \fuente{TP2, p. 1} \\
\bottomrule
\end{tabularx}
\medskip

\subsection{Cifras para tener a mano}
\begin{tabularx}{\linewidth}{@{}LP{4.2cm}P{2.6cm}@{}}
\toprule
\enc{Qué} & \enc{Valor} & \enc{Fuente} \\
\midrule
Filas × columnas & 41.188 × 21 & \fuente{CSV} \\
«yes» / «no» & 4.640 (11,27\%) / 36.548 & \fuente{CSV} \\
Exactitud de decir siempre «no» & 88,73\% & \fuente{CSV} \\
AUC de \texttt{duration} sola & 0,818 & \fuente{CSV} \\
Filas con algún \texttt{unknown} & 10.700 (25,98\%) & \fuente{CSV} \\
\texttt{pdays = 999} / con contacto previo & 39.673 (96,32\%) / 4.110 & \fuente{CSV} \\
«yes» por año inferido & 4,84\% · 19,48\% · 52,14\% & \fuente{names §4 · CSV} \\
Correlación del bloque macro & 0,91 a 0,97 & \fuente{CSV} \\
Varianza de \texttt{duration} + \texttt{pdays} & 95,0\% del total & \fuente{CSV} \\
Duplicados exactos & 12 & \fuente{CSV} \\
\bottomrule
\end{tabularx}

\anexo{Glosario}
\begin{glosario}
\termino{Depósito a plazo fijo}{\textit{Term deposit}: el producto que se ofrece en la llamada; $y$ dice si el cliente lo contrató.}
\termino{Fuga de datos}{\textit{Data leakage}: información que debería estar sólo del lado de la evaluación entra al entrenamiento, y las métricas mienten hacia arriba. En el TP2, también una variable que no existe antes de llamar.}
\termino{Centinela}{Un valor que no es un número sino una marca, como el 999 de \texttt{pdays}.}
\termino{Curva de validación}{\textit{Validation curve}: cómo varía el rendimiento al cambiar un hiperparámetro \fuente{Clase 7, slide 50}; el TP2 pide usarlas «para elegir el mejor valor» \fuente{TP2, p. 2}.}
\end{glosario}

\end{document}
